% © 2026 Ing. David Jaroš — CC BY-NC-ND 4.0
% Status: audit note; GAP-10T-DYN and GAP-10D remain NARROWED (see verdict).
\documentclass[11pt,a4paper]{article}
\usepackage{amsmath,amssymb,amsthm}
\usepackage[most]{tcolorbox}
\usepackage{booktabs}
\usepackage[margin=2.5cm]{geometry}
\newtheorem{theorem}{Theorem}
\newtheorem{lemma}{Lemma}
\newtheorem{corollary}{Corollary}
\newcommand{\B}{\mathbb B}
\newcommand{\dd}{\ddagger}
\title{Canonical UBT Action Audit for GAP-10T-DYN and GAP-10D:\\
Dependency Classification, Exact Spin Current, and a Flat Affine No-Go}
\author{Ing. David Jaro\v{s}}
\date{26 July 2026}
\begin{document}\maketitle

\begin{abstract}
We audit whether the canonical UBT action, as currently written in the
repository, derives the Cartan torsion equation and the
Einstein--$\Lambda$ tetrad equation without importing them.  The answer is
negative and is now made precise.  We classify every term of the working
action into derived, assumed, imported, or undefined.  We derive, exactly,
the tree-level spin current of the canonical kinetic term for the pure-pair
representative $A_\mu=\Omega_\mu$, $B_\mu=-\Omega_\mu^{\dd}$, prove that it
couples only to the anti-Hermitian part of $\Theta$ whenever
$D_\mu\Theta$ lies in the Lorentz slice, and prove a pointwise rigidity
statement: for a nondegenerate tetrad the current vanishes iff the
anti-Hermitian part of $\Theta$ vanishes.  Consequently the minimal
Hilbert--Palatini $+$ kinetic branch forces nonzero torsion on the
complement of at most one point of \emph{every} flat affine representer,
so the flat inertial solution of GAP-10I-SR is not a torsion-free solution
of the minimal branch.  This is an exact flat-space no-go that corroborates,
from an independent direction, the curved Gaussian-patch scaling obstruction.
GAP-10T-DYN and GAP-10D remain NARROWED with the remaining lemmas named
explicitly.  All algebraic statements are verified by
\texttt{tools/verify\_canonical\_spin\_current.py}.
\end{abstract}

\section{Canonical field content}
The single fundamental field is $\Theta(q,\tau)\in\B\simeq\mathrm{Mat}(2,\mathbb C)$.
The involution $\dd$ is Hermitian conjugation; $\sharp$ is the $2\times2$
adjugate, $X^\sharp=\operatorname{tr}(X)\mathbf1-X$.  The classical Lorentz
slice is $W_L=\{X:\,X^{\dd}=-X\}$ with basis $E_0=i\mathbf1$,
$E_k=-i\sigma_k$, and the central Jordan identity
$\tfrac12(X^\sharp Y+Y^\sharp X)=\eta(X,Y)\mathbf1$ holds on $W_L$
[verified: C1].  The covariant jet is
$E_\mu=\mathcal N_0^{-1/2}D_\mu\Theta$ with
\begin{equation}
  D_\mu\Theta=\partial_\mu\Theta+\Omega_\mu\Theta+\Theta\,\Omega_\mu^{\dd},
  \qquad \Omega_\mu\in\mathfrak{sl}(2,\mathbb C),
  \label{eq:purepair}
\end{equation}
which is the unique no-new-field representative established by
GAP-10I-PAIR-KIN.  The connection action preserves $W_L$ [C1].

\section{Existing action}
The working action in \texttt{canonical/THEORY/canonical/canonical\_action.tex}
is
\begin{equation}
  S \;=\; S_{\rm HP}[e,\omega] \;+\; S_\Theta,\qquad
  S_\Theta=\tfrac12\int\sqrt{-g}\,\big\langle D_\mu\Theta,\,D^\mu\Theta\big\rangle
  -\kappa V[\Theta],
  \label{eq:action}
\end{equation}
with $S_{\rm HP}$ the Hilbert--Palatini/Einstein--Cartan first-order term
with cosmological constant.  Two pairings occur in the corpus and both are
audited throughout: the $\dd$ Hilbert--Schmidt pairing
$\langle X,Y\rangle_{\dd}=\operatorname{Re}\operatorname{tr}(X^{\dd}Y)$
and the $\sharp$ scalar pairing
$\langle X,Y\rangle_{\sharp}=\operatorname{Re}\operatorname{Sc}(X^\sharp Y)$.

\subsection*{Dependency classification (audit item 2)}
\begin{center}
\begin{tabular}{@{}llp{7.2cm}@{}}
\toprule
Object & Class & Basis \\
\midrule
$\Theta$, $\B$, $\dd$, $\sharp$, $W_L$ & (a) derived/definitional & core algebra files \\
$E_\mu=\mathcal N_0^{-1/2}D_\mu\Theta$ & (a) derived & covariant-tetrad chain \\
$g_{\mu\nu}$ from central Jordan identity & (a) derived & rank theorem [L1] \\
$\Omega_\mu$ pure-pair form \eqref{eq:purepair} & (a) derived & GAP-10I-PAIR-KIN [L1] \\
$S_\Theta$ kinetic form and pairing choice & (b) low-energy assumption & form assumed; pairing not selected canonically \\
Spin current $\tau^{\mu}{}_{ab}$ & (a) \textbf{derived here} & Theorem~\ref{thm:spincurrent} \\
$S_{\rm HP}$, $\Lambda$ term & (c) standard import & stated as effective branch \\
$\kappa$, sign and value & (d) undefined from UBT & \S\,Coefficient matching \\
Matter stress tensor $T_{\mu\nu}$ & (a) derived given \eqref{eq:action} & \texttt{canonical/t\_munu/} \\
Potential $V[\Theta]$ & (b) assumption & unconstrained by this audit \\
\bottomrule
\end{tabular}
\end{center}

\section{Independent versus composite variations}
Two inequivalent schemes exist and must not be conflated.
\emph{Independent (Palatini) scheme:} $e$ and $\omega$ vary independently;
$S_\Theta$ depends on $\omega$ only algebraically --- the kinetic density is a
polynomial of total degree two in the components of $\Omega_\mu$ with no
derivatives of $\Omega_\mu$ [C2].  Hence the $\omega$-variation of
$S_\Theta$ is an algebraic current and \emph{no curvature invariant arises
from $S_\Theta$ at tree level}; every curvature term in the field equations
originates in the imported $S_{\rm HP}$.
\emph{Composite scheme:} substituting $\omega=\mathring\omega(e)+K$ before
variation makes $S_\Theta$ depend on $\partial e$; the induced quadratic form
in $\partial e$ is $\Theta$-dependent and is \emph{not} shown to reproduce
the Einstein--Hilbert contraction.  Deciding this is part of the named
remaining lemma below, not a result of this audit.

\section{Torsion equation}
In the independent scheme, $\delta S/\delta\omega$ yields the Cartan equation
$\epsilon_{abcd}T^a\wedge e^b=\kappa\,\tau_{cd}$ with the $24\times24$ torsion
map pointwise invertible for a nondegenerate tetrad (GAP-10T-PALATINI [L1]).
The new content is the source itself.

\begin{theorem}[Exact tree-level spin current]\label{thm:spincurrent}
For the pure-pair representative \eqref{eq:purepair} and either pairing, the
variation $\Omega_\mu\to\Omega_\mu+\varepsilon M$, $M\in\mathfrak{sl}(2,\mathbb C)$,
gives
\begin{equation}
  \tau^{\mu}(M)\;=\;\big\langle\,M\Theta+\Theta M^{\dd},\;D^\mu\Theta\,\big\rangle .
  \label{eq:tau}
\end{equation}
\end{theorem}

\begin{lemma}[Slice lemma]\label{lem:slice}
If $D^\mu\Theta\in W_L$ (as required by the tetrad construction), then
$\tau^\mu(M)$ depends only on the anti-Hermitian part of $\Theta$; the
Hermitian part decouples identically.  Verified exactly for both pairings
[C3].
\end{lemma}

\begin{theorem}[Pointwise rigidity]\label{thm:rigidity}
Let $D_\mu\Theta=\sqrt{\mathcal N_0}\,E_\mu$ with $\{E_\mu\}$ the standard
nondegenerate slice basis.  Then $\tau^\mu(M)=0$ for all $\mu$ and all
$M\in\mathfrak{sl}(2,\mathbb C)$ iff the anti-Hermitian part of $\Theta$
vanishes at that point.  Verified exactly for both pairings [C4].
\end{theorem}

\begin{corollary}[Flat affine no-go]\label{cor:nogo}
On every affine representer
$\Theta=\Theta_0+\sqrt{\mathcal N_0}\,E_\nu x^\nu$ (GAP-10I-SR), the
anti-Hermitian part of $\Theta(x)$ is an affine non-constant $W_L$-valued
map and vanishes at most at one point.  Hence $\tau\neq0$ on the complement
of at most one point; the $x$-gradient of every component of $\tau$ is a
$\Theta_0$-independent constant ($\pm2\mathcal N_0$ for $\dd$,
$\pm\mathcal N_0$ for $\sharp$) [C5].  By invertibility of the Cartan map,
the minimal branch \eqref{eq:action} forces nonzero torsion there, so the
flat inertial, torsion-free solution is \emph{not} a solution of the
minimal branch.
\end{corollary}

This is an exact flat-space counterpart of the curved Gaussian-patch scaling
obstruction ($K\sim1/\rho$ kinematic versus $\tau\sim\rho$ sourced): both
identify the same structural mismatch of the minimal branch, by independent
arguments.

\section{Tetrad equation}
In the independent scheme, $\delta S/\delta e$ yields
$G_{\mu\nu}+\Lambda g_{\mu\nu}=\kappa T_{\mu\nu}$ with $T_{\mu\nu}$ the
canonical stress tensor of $S_\Theta$; the Einstein tensor originates
entirely in $S_{\rm HP}$ [C2].  No derivation of the tetrad equation from
$S_\Theta$ alone exists in the corpus, and this audit proves none.

\section{Coefficient matching}
Neither the sign nor the magnitude of $1/4\kappa$, nor $\Lambda$, is fixed
by $\mathcal N_0$, by the $\Theta$ vacuum, or by any repository derivation.
The only induced-gravity infrastructure present
(\texttt{scripts/spectral/heat\_kernel\_trace.py}) computes flat-torus
spectral traces and does not reach the curved Seeley--DeWitt $a_2$
coefficient.  The coefficient therefore remains an explicit open subgap
(part of GAP-10D).

\section{No-circularity audit}
(i) The spin current \eqref{eq:tau} is derived from $S_\Theta$ alone; no
Einstein or Cartan equation is substituted anywhere in
Theorems~\ref{thm:spincurrent}--\ref{thm:rigidity}.
(ii) The Cartan equation is used only \emph{after} the current is derived,
and only its proved algebraic invertibility [L1] is invoked.
(iii) No step assumes the conclusion of GAP-10D; the Einstein--$\Lambda$
endpoint is cited exclusively as a property of the imported $S_{\rm HP}$.
(iv) No new independent fields are introduced; the audit works entirely in
the pure-pair representative.

\section{Exact closure verdict}
\begin{tcolorbox}[colback=orange!6!white,colframe=orange!70!black,
title=Verdict]
\textbf{GAP-10T-DYN: NARROWED} (not closed, not closable in the minimal
branch as written).  Newly derived inside it:
the exact tree-level spin current (Theorem~\ref{thm:spincurrent},
conditional [L1] on the stated pairing class and the pure-pair
representative) and the flat affine no-go
(Corollary~\ref{cor:nogo}, [L1]).  \textbf{Remaining lemma, named:}
construct a canonical modification --- a non-minimal torsion term, a
canonically selected pairing/projection annihilating the affine current, or
a torsion-free relative bimodule completion --- whose torsion sector admits
the flat affine representer, and derive the selected branch from
$S[\Theta]$.

\medskip
\textbf{GAP-10D: NARROWED.}  The Einstein--$\Lambda$ endpoint remains
conditional on the imported $S_{\rm HP}$ and Lovelock hypotheses
(GAP-10D-PALATINI / GAP-10D-UNIQUENESS, unchanged).  \textbf{Remaining
lemma, named (induced-Palatini lemma):} show that the $\Theta$ effective
action contains $\epsilon_{abcd}e^a\wedge e^b\wedge R^{cd}(\omega)$ with a
sign and coefficient fixed by $(\mathcal N_0,\Theta$-vacuum$)$, e.g.\ via
the curved heat-kernel $a_2$ coefficient of $D^\dagger D$; the present
spectral tooling stops at flat tori.

\medskip
Prohibited conclusions, explicitly not claimed: closure of GAP-10T-DYN or
GAP-10D; derivation of $S_{\rm HP}$; any statement that the composite
scheme reproduces Einstein--Hilbert.
\end{tcolorbox}

\end{document}
